% © 2026 Ing. David Jaroš — CC BY-NC-ND 4.0
%
% This work is licensed under a Creative Commons Attribution-NonCommercial-NoDerivatives
% 4.0 International License (CC BY-NC-ND 4.0).
%
% License History: Earlier drafts (up to v0.3) were released under CC BY 4.0.
% From v0.4 onward, all material is released under CC BY-NC-ND 4.0 to protect
% the integrity of the theoretical work during ongoing academic development.
%
% See LICENSE.md for full license text.

% Mirror Sector: Vacuum Stability Analysis
% Status: [NEW CONJECTURE — motivated by twin prime observation]
% Date: 2026-03-07

\ifdefined\INCLUDEMODE
\else
  \documentclass[12pt]{article}
  \usepackage[a4paper, margin=2.5cm]{geometry}
  \usepackage{amsmath, amssymb, amsthm}
  \usepackage{hyperref}
  \usepackage{xcolor}
  \usepackage{mdframed}
  \usepackage{booktabs}
  \usepackage{tcolorbox}

  \newtheorem{theorem}{Theorem}[section]
  \newtheorem{lemma}[theorem]{Lemma}
  \newtheorem{proposition}[theorem]{Proposition}
  \theoremstyle{definition}
  \newtheorem{definition}[theorem]{Definition}
  \theoremstyle{remark}
  \newtheorem{remark}[theorem]{Remark}

  \newmdenv[backgroundcolor=yellow!20, linecolor=orange, linewidth=1.5pt]{gapbox}
  \newmdenv[backgroundcolor=green!15, linecolor=green!60!black, linewidth=1.5pt]{resultbox}
  \newmdenv[backgroundcolor=blue!10, linecolor=blue!60, linewidth=1.5pt]{conjecturebox}

  \tcbuselibrary{skins,breakable}
  \newtcolorbox{statusbox}[1]{
    colback=yellow!8!white, colframe=orange!60!black,
    title={\textbf{#1}}, fonttitle=\bfseries, breakable
  }

  \title{\textbf{Mirror Sector: Vacuum Stability Analysis}\\
         \large $n^* = 139$ as a Metastable Vacuum in UBT}
  \author{David Jaroš}
  \date{2026-03-07}

  \begin{document}
  \maketitle

  \begin{abstract}
  The twin prime observation in the UBT Hecke analysis ($n^*=137$ global vacuum,
  $n^*=139$ local minimum) motivates a study of vacuum stability in the effective
  potential $V_{\rm eff}(n) = A\cdot n^2 - B\cdot n\ln(n)$.
  We compute $V_{\rm eff}(137)$ and $V_{\rm eff}(139)$ numerically, estimate the
  Coleman tunnelling rate $\Gamma \sim e^{-S_E}$, and compute the corresponding
  vacuum lifetime $\tau$.  The prediction $\tau \gg t_{\rm universe}$ (required by
  the anthropic principle) is checked.

  \textbf{Status: [NEW CONJECTURE — motivated by twin prime observation]}
  \end{abstract}
\fi

\section{Mirror Sector: Vacuum Stability}
\label{sec:vacuum_stability}

\begin{statusbox}{Status: [NEW CONJECTURE — motivated by twin prime observation]}
This analysis is motivated by the numerical observation that Hecke eigenvalues at
primes $p=137$ (Set A, our sector) and $p=139$ (Set B, mirror sector) are both
globally unique in the range $50$--$300$.  This suggests that $n^*=137$ is the
global vacuum and $n^*=139$ is a metastable local minimum of the UBT effective
potential.  All conclusions are conjectural at this stage.
\end{statusbox}

\subsection{The Effective Potential}

The UBT effective potential for the vacuum mode number $n$ takes the one-loop form:
\begin{equation}
  V_{\rm eff}(n) = A\cdot n^2 - B\cdot n\ln(n),
  \label{eq:Veff}
\end{equation}
where:
\begin{itemize}
  \item $A$ is the quadratic coefficient from the kinetic term (positive; sets energy scale),
  \item $B$ is the logarithmic correction from one-loop running (positive; destabilises).
\end{itemize}

The minimum of $V_{\rm eff}(n)$ at $n = n^*$ satisfies $V_{\rm eff}'(n^*) = 0$:
\begin{equation}
  2A\cdot n^* = B\cdot (\ln n^* + 1)
  \;\Longrightarrow\;
  n^* = \frac{B}{2A}\left(\ln n^* + 1\right).
  \label{eq:minimum_condition}
\end{equation}
The observed value $n^* = 137$ gives the calibration relation
$B_0 = 2A\cdot 137 / (\ln 137 + 1)$.

\subsection{Numerical Evaluation: $V_{\rm eff}(137)$ vs $V_{\rm eff}(139)$}

We evaluate $V_{\rm eff}$ at the two twin prime values.  Define the dimensionless
ratio $r = B/(2A)$.  From the calibration to $n^* = 137$:
\begin{equation}
  r = \frac{n^*}{\ln n^* + 1} = \frac{137}{\ln 137 + 1} \approx \frac{137}{5.919 + 1} = \frac{137}{6.919} \approx 19.80.
  \label{eq:r_value}
\end{equation}

With $A = 1$ (setting energy units), $B = 2r \approx 39.59$:

\subsubsection{At $n = 137$:}
\begin{align}
  V_{\rm eff}(137) &= 137^2 - 39.59 \cdot 137 \cdot \ln(137) \notag\\
  &= 18769 - 39.59 \times 137 \times 4.920 \notag\\
  &= 18769 - 39.59 \times 674.04 \notag\\
  &= 18769 - 26684.4 \notag\\
  &\approx -7915.
  \label{eq:Veff137}
\end{align}

\subsubsection{At $n = 139$:}
\begin{align}
  V_{\rm eff}(139) &= 139^2 - 39.59 \cdot 139 \cdot \ln(139) \notag\\
  &= 19321 - 39.59 \times 139 \times 4.934 \notag\\
  &= 19321 - 39.59 \times 686.03 \notag\\
  &= 19321 - 27158.5 \notag\\
  &\approx -7838.
  \label{eq:Veff139}
\end{align}

\begin{resultbox}
\textbf{[NUMERICAL OBSERVATION]:}
$V_{\rm eff}(137) \approx -7915 < V_{\rm eff}(139) \approx -7838$.

The global minimum is at $n^* = 137$ (our sector).
The value at $n = 139$ is higher by $\Delta V \approx 77$ (in units where $A=1$),
consistent with $n^* = 139$ being a local but not global minimum --- a metastable
vacuum.
\end{resultbox}

\subsubsection{Second Derivative Test}

To confirm $n = 139$ is a local minimum (metastable), we check $V_{\rm eff}''(139) > 0$:
\begin{equation}
  V_{\rm eff}''(n) = 2A - \frac{B}{n} = 2 - \frac{39.59}{n}.
\end{equation}
At $n = 139$: $V_{\rm eff}''(139) = 2 - 39.59/139 \approx 2 - 0.285 = 1.715 > 0$. ✓

At $n = 137$: $V_{\rm eff}''(137) = 2 - 39.59/137 \approx 2 - 0.289 = 1.711 > 0$. ✓

Both are local minima; $n^* = 137$ is deeper (global).

\subsection{Coleman Tunnelling Estimate}

The metastable vacuum $n^* = 139$ can tunnel to the true vacuum $n^* = 137$ via
Coleman's bounce solution \cite{Coleman1977}.  The Euclidean action for the bounce is
approximately:
\begin{equation}
  S_E \approx \frac{\pi^2 \cdot \Delta V \cdot (\Delta n)^4}{\sigma^3},
  \label{eq:coleman_action}
\end{equation}
where:
\begin{itemize}
  \item $\Delta V = V_{\rm eff}(139) - V_{\rm eff}(137) \approx 77$ (energy difference),
  \item $\Delta n = 139 - 137 = 2$ (width of the barrier in $n$-space),
  \item $\sigma$ is the effective surface tension of the domain wall
        (estimated as $\sigma \sim \sqrt{A \cdot \Delta V} \cdot \Delta n$).
\end{itemize}

Estimating the surface tension:
\begin{equation}
  \sigma \sim \sqrt{\Delta V} \cdot \Delta n = \sqrt{77} \times 2 \approx 8.77 \times 2 = 17.5.
\end{equation}

Euclidean action:
\begin{equation}
  S_E \approx \frac{\pi^2 \times 77 \times 16}{17.5^3}
  = \frac{12168}{5359} \approx 2.27.
  \label{eq:SE_value}
\end{equation}

\begin{gapbox}
\textbf{Note}: The surface tension estimate $\sigma \sim \sqrt{A\cdot\Delta V}\cdot\Delta n$
is a rough approximation using the thin-wall formula.  The actual $\sigma$ requires
knowledge of the potential profile between $n=137$ and $n=139$, which depends on the
continuous extension of $V_{\rm eff}$ to non-integer $n$.  The estimate above should
be treated as an order-of-magnitude result only.
\end{gapbox}

\subsection{Vacuum Lifetime}

The tunnelling rate per unit volume is:
\begin{equation}
  \Gamma/V \sim \Lambda^4 \cdot e^{-S_E},
\end{equation}
where $\Lambda$ is the relevant energy scale (the compactification scale,
$\Lambda \sim 1/R_\psi$).  The total decay rate for the observable universe
(volume $V_{\rm universe} \sim (c/H_0)^3$) is:
\begin{equation}
  \Gamma_{\rm total} = (\Gamma/V) \times V_{\rm universe}.
\end{equation}
The vacuum lifetime is:
\begin{equation}
  \tau = 1/\Gamma_{\rm total} \sim \frac{e^{S_E}}{\Lambda^4 V_{\rm universe}}.
  \label{eq:lifetime}
\end{equation}

With $S_E \approx 2.27$ and $\Lambda \sim m_e c^2 / \hbar \sim 10^{21}$ Hz,
$V_{\rm universe} \sim (10^{26} \text{ m})^3$, we get:
\begin{equation}
  \tau \sim e^{2.27} \times \frac{1}{\Lambda^4 V_{\rm universe}}
  \sim 9.7 \times \frac{\hbar^4}{\Lambda^4 \hbar^3 c^3} \times \frac{1}{(c/H_0)^3}.
\end{equation}

\begin{gapbox}
\textbf{Gap}: The Euclidean action $S_E \approx 2.27$ is very small, implying a
short lifetime $\tau \ll t_{\rm universe}$ in this rough estimate.  This is
\emph{inconsistent} with the anthropic requirement $\tau > t_{\rm universe}$.

\medskip
\textbf{Resolution candidates}:
\begin{enumerate}
  \item The $n$-variable may be discrete (integer-valued), in which case the
        Coleman formula for continuous tunnelling does not directly apply.
        A discrete tunnelling rate would be exponentially suppressed differently.
  \item The surface tension $\sigma$ may be much larger if the barrier between
        $n=137$ and $n=139$ involves multiple minima or a large intermediate potential.
  \item The physical interpretation of $n$ as a vacuum mode number (not a spatial
        field) means the ``volume'' $V_{\rm universe}$ should not appear in the
        naive Coleman formula; the relevant ``volume'' may be $\sim R_\psi^3 \ll V_{\rm universe}$.
\end{enumerate}
\textbf{Status}: Coleman estimate is inconclusive; a careful treatment of the
discrete/compact nature of the $n$-variable is required.
\end{gapbox}

\subsection{Prediction: Anthropic Constraint}

For consistency with our existence:
\begin{equation}
  \tau > t_{\rm universe} \approx 4.35 \times 10^{17}~\text{s}.
  \label{eq:anthropic}
\end{equation}
This requires $S_E \gg \ln(t_{\rm universe} \Lambda^4 V_{\rm universe}/\hbar)$,
which is a very large number $\sim 10^3$ in standard cosmological units.

\begin{conjecturebox}
\textbf{[MOTIVATED CONJECTURE]:}
$n^* = 139$ is a metastable vacuum of $V_{\rm eff}(n)$ with a lifetime
$\tau \gg t_{\rm universe}$.  The numerical estimate above with the naive
Coleman formula gives $S_E \approx 2.27$ (too small), suggesting that either:
(a) the correct treatment of discrete-$n$ tunnelling gives a much larger $S_E$, or
(b) the volume factor in the Coleman formula should be $R_\psi^3$ rather than
$V_{\rm universe}$, which would push $\tau$ to anthropically safe values.

\medskip
This conjecture is required for consistency of the mirror sector interpretation:
if the mirror vacuum $n^*=139$ were short-lived, the mirror sector could not
support stable matter structures.
\end{conjecturebox}

\subsection{Summary}

\begin{center}
\renewcommand{\arraystretch}{1.3}
\begin{tabular}{lll}
\toprule
\textbf{Result} & \textbf{Status} & \textbf{Notes} \\
\midrule
$V_{\rm eff}(137) < V_{\rm eff}(139)$ & \textbf{[NUMERICAL OBSERVATION]} &
  Follows from $V_{\rm eff}$ form with $B > B_0$ \\
$n^* = 139$ is local minimum & \textbf{[NUMERICAL OBSERVATION]} &
  $V_{\rm eff}''(139) > 0$ confirmed \\
$n^* = 139$ is metastable vacuum & \textbf{[MOTIVATED CONJECTURE]} &
  Coleman tunnelling; $\tau$ estimate inconclusive \\
$\tau > t_{\rm universe}$ & \textbf{[CONJECTURE]} &
  Required by anthropics; naive estimate fails \\
\bottomrule
\end{tabular}
\end{center}

\noindent\textbf{Date}: 2026-03-07 \\
\textbf{Author}: David Jaroš

\begin{thebibliography}{9}
\bibitem{Coleman1977}
S.\ Coleman, \textit{Fate of the false vacuum: Semiclassical theory},
Phys.\ Rev.\ D \textbf{15}, 2929 (1977).

\bibitem{MirrorSector}
R.\ Foot and R.\ R.\ Volkas, \textit{Neutrino physics and the mirror world:
How exact parity symmetry explains the solar neutrino deficit, the atmospheric
neutrino anomaly and the LSND experiment},
Phys.\ Rev.\ D \textbf{52}, 6595 (1995).
\end{thebibliography}

\ifdefined\INCLUDEMODE
\else
\end{document}
\fi
